% ============================================================
\chapter{Resultados y evaluación}
\label{cap:resultados}
% ============================================================

Este capítulo presenta y analiza los resultados obtenidos en cada una de las tres
capas del sistema \textit{kubescan}, seguidos de una validación cruzada con
herramientas de análisis estático externas y de una evaluación funcional de la
herramienta como software distribuible. Los resultados se organizan en el mismo
orden que las fases de entrenamiento descritas en el Capítulo~\ref{cap:diseno}. Las
métricas empleadas se justifican en detalle en la Sección~\ref{sec:metricas}.

% -----------------------------------------------------------
\section{Resultados de la Capa 1: Random Forest}
% -----------------------------------------------------------

\subsection{Clasificación binaria}

La Tabla~\ref{tab:rf_results} resume los resultados del modelo binario mediante
cuatro métricas: F1 macro (media no ponderada de F1 entre las clases seguro y
misconfigured, que da igual peso a ambas independientemente de su frecuencia en el
dataset), exactitud (\textit{accuracy}, porcentaje simple de aciertos), \ac{AUC}-\ac{ROC}
(capacidad del modelo para separar las dos clases a cualquier umbral de decisión) y
\textit{out-of-bag score} (estimación interna de generalización que
\textit{Random Forest} calcula sobre las muestras que cada árbol no vio durante su
entrenamiento, sin necesitar un conjunto de validación aparte). El clasificador
alcanza una F1 macro de 0,9946 en el conjunto de test, superando el objetivo de
diseño (OE2, F1 macro\,>\,0,85) en 14,5 puntos porcentuales.

\begin{table}[htbp]
\centering
\caption{Resultados del Random Forest — clasificación binaria}
\label{tab:rf_results}
\begin{tabular}{@{}lrrr@{}}
\toprule
Métrica & Test & \acs{CV} (5-fold) & Objetivo \\
\midrule
F1 macro & \textbf{0,9946} & 0,9917 \(\pm\) 0,0053 & \(>0{,}85\) \checkmark \\
Exactitud (\textit{accuracy}) & 0,9947 & 0,9919 \(\pm\) 0,0052 & — \\
\ac{AUC}-\ac{ROC} & \textbf{0,9986} & 0,9984 \(\pm\) 0,0011 & — \\
\textit{Out-of-bag score} & 0,9920 & — & — \\
\bottomrule
\end{tabular}
\fuentepropia
\end{table}

La Figura~\ref{fig:rf_confusion} muestra la matriz de confusión sobre el conjunto
de test (566~muestras). El modelo comete sólo 3~errores totales: 3~falsos positivos
(manifiestos seguros clasificados como misconfigured) y 0~falsos negativos.
La ausencia de falsos negativos es la propiedad más crítica para seguridad operacional:
significa que ningún manifiesto con misconfiguraciones reales escapa al filtrado de
la Capa~1, independientemente de cuántas flags estén activas o de su rareza en el
dataset. Los 3~falsos positivos corresponden a manifiestos seguros que acumulan varias
features de riesgo de severidad media (por ejemplo, \texttt{NO\_RUN\_AS\_NON\_ROOT}
y \texttt{IMAGE\_USES\_LATEST} simultáneamente), pero que no tienen ninguna
flag de escape activa; el operador de seguridad puede descartarlos rápidamente
inspeccionando el campo \texttt{flags} de la salida \ac{JSON}.

El \ac{AUC}-\ac{ROC} de 0,9986 (prácticamente~1,00) indica que el modelo puede ordenar las
muestras por probabilidad de misconfiguration de forma casi perfecta a cualquier umbral
de decisión. En la práctica, este valor implica que, si se ordenan los 566~manifiestos
de test de mayor a menor \textit{risk\_score}, los manifiestos misconfigured aparecen
casi todos antes que los seguros. Esta propiedad es la que habilita el uso del
\textit{risk\_score} como característica nodal cuantitativa en el grafo de la Capa~2:
la señal transmitida a la \ac{GNN} no es un bit binario sino una probabilidad calibrada
que refleja la densidad de misconfiguraciones del manifiesto.

\begin{figure}[htbp]
  \centering
  \caption{Matriz de confusión del Random Forest en el conjunto de test}
  \label{fig:rf_confusion}
  \includegraphics[width=0.52\textwidth]{fig_rf_confusion}

  \fuentepropia
\end{figure}

En la matriz de la Figura~\ref{fig:rf_confusion} (566~muestras), los 3~falsos
positivos son manifiestos seguros con alta densidad de características de riesgo;
no existe ningún falso negativo.

\subsection{Importancia de características}
\label{sec:rf_importance}

La Figura~\ref{fig:rf_importance} muestra las 10 características más relevantes del
modelo. La característica más discriminativa es \texttt{total\_misconfigs}
(33,05\,\% de importancia), seguida de \texttt{NO\_RESO} (14,79\,\%) e
\texttt{INSECURE\_HTTP} (11,90\,\%). Las flags de escape críticas (\texttt{TRUE\_HOST\_PID},
\texttt{SEC\_CONT\_OVER\_PRIVIL}) presentan importancias menores (\(<\)1\,\%) por su
rareza en el dataset, pero son identificadas correctamente cuando están presentes,
como confirma la ausencia de falsos negativos en la clasificación binaria.

\begin{figure}[htbp]
  \centering
  \caption{Importancia de características del clasificador Random Forest}
  \label{fig:rf_importance}
  \includegraphics[width=\textwidth]{fig_rf_importance}

  \fuentepropia
\end{figure}

Las dos primeras características de la Figura~\ref{fig:rf_importance} (en rojo)
concentran el 48\,\% de la capacidad discriminativa total del modelo.

Un estudio de ablación eliminando \texttt{total\_misconfigs} produce F1 idéntico
y la misma matriz de confusión, confirmando que el modelo no depende de esta
característica derivada y que no existe fuga de datos circular.

\subsection{Modelo de severidad}

El clasificador de tres clases (\textit{clean}, \textit{low-medium},
\textit{high-critical}) alcanza F1 macro\,=\,0,9838 en test. La Tabla~\ref{tab:severity_f1}
desglosa por clase la precisión (proporción de predicciones de esa clase que son
correctas), el \textit{recall} (proporción de casos reales de esa clase que el
modelo identifica), el F1 (media armónica de ambas) y el soporte (número de
muestras de test en esa clase):

\begin{table}[htbp]
\centering
\caption{F1 por clase — modelo de severidad (3 clases, test)}
\label{tab:severity_f1}
\begin{tabular}{@{}lrrrr@{}}
\toprule
Clase & Precisión & Recall & F1 & Soporte \\
\midrule
\textit{clean}          & 1,0000 & 0,9906 & 0,9953 & 319 \\
\textit{low-medium}     & 0,9615 & 0,9934 & 0,9772 & 151 \\
\textit{high-critical}  & 0,9894 & 0,9688 & 0,9789 & 96 \\
\midrule
Macro media & 0,9836 & 0,9842 & 0,9838 & — \\
\bottomrule
\end{tabular}
\fuentepropia
\end{table}

La clase \textit{high-critical} (96~muestras de test, que representan manifiestos con
flags de escape activas como \texttt{TRUE\_HOST\_PID} o \texttt{SEC\_CONT\_OVER\_PRIVIL})
obtiene el F1 más bajo de las tres clases (0,9789), con 3~errores, los tres clasificados
como \textit{low-medium}; ningún manifiesto crítico es clasificado como \textit{clean}.
Esta ausencia de fugas hacia la clase \textit{clean} es la propiedad más relevante desde
el punto de vista de seguridad operacional: en el peor caso, un manifiesto crítico es
subestimado a \textit{low-medium}, nunca descartado como seguro. El resto de errores del
modelo se concentra en la frontera entre \textit{clean} y \textit{low-medium} (3~casos), donde
algunos manifiestos con pocas flags de baja severidad se sobreclasifican como
\textit{low-medium} (falsos positivos benignos desde el punto de vista del operador); un
único manifiesto \textit{low-medium} se sobreclasifica como \textit{high-critical}, un error
igualmente conservador.

El modelo de severidad se utiliza en el informe de salida de \textit{kubescan} para
proporcionar al operador una descripción cualitativa del nivel de riesgo de cada
manifiesto, complementando el \textit{risk\_score} numérico del modelo binario.

% -----------------------------------------------------------
\section{Resultados de la Capa 2: \acs{GNN}}
\label{sec:layer2}
% -----------------------------------------------------------

\subsection{Evolución por fase experimental}

La Tabla~\ref{tab:gnn_evolution} muestra la evolución de los resultados de la \ac{GAT}
a lo largo de las seis fases del desarrollo experimental. La
Figura~\ref{fig:gnn_evolution} traza la trayectoria de las cuatro
primeras fases y muestra cómo la aumentación de datos y la corrección del vector
de escape \texttt{HOSTPATH\_MOUNT} contribuyeron de forma independiente y
complementaria. La quinta fase corrige el protocolo de particionado: las fases~1--4
incluían en el entrenamiento variantes aumentadas derivadas de los
clústeres de validación y test (fuga de datos por casi-duplicados), por lo que sus
métricas están infladas de forma optimista. Esa quinta fase aplica el protocolo
\textit{group-aware} descrito en el Capítulo~\ref{cap:diseno} (las variantes acompañan a su
clúster base y los clústeres de test quedan fuera de los pliegues) e incorpora el
\textit{embedding} de tipos de arista; sus métricas son metodológicamente válidas
pero corresponden todavía al corpus preliminar de 471~grafos. La sexta y última
fase incorpora el corpus completo de Rahman et~al.\ (599~grafos originales,
37~cadenas; Sección~\ref{sec:construccion_dataset}) y extiende el protocolo de
particionado a las \textbf{familias de plantilla} (Capítulo~\ref{cap:diseno}):
durante esta fase se detectó que los casi duplicados de una misma plantilla
(ocho \texttt{badpods\_*}, nueve \texttt{datadog\_*}, ocho
\texttt{simulator\_*}, doce \texttt{kubernetes\_goat\_*}) quedaban repartidos
entre particiones, una fuga de información adicional a la de las variantes
aumentadas que inflaba las métricas de las fases 1--5. Con las familias atómicas,
entropía cruzada ponderada como función de pérdida y selección de
\textit{checkpoints} por Precisión@5 (ablación en la
Sección~\ref{sec:ablacion_perdida}), esta fase es la que se reporta como
resultado final del trabajo — sus métricas son las que produce el modelo
efectivamente desplegado y reproducible a partir de los \textit{checkpoints} del
repositorio. Por el cambio de protocolo, sus cifras no son directamente
comparables con las de las fases anteriores: miden generalización entre familias,
no interpolación entre casi duplicados.

\begin{table}[htbp]
\centering\small
\caption{Evolución de resultados de la Capa~2 (\ac{GAT}, 5-fold \ac{CV})}
\label{tab:gnn_evolution}
\begin{tabular}{@{}p{4.3cm}rrrrc@{}}
\toprule
Fase & Grafos & Cadenas & F1 macro & P@5 & \shortstack{Cobert.\\techo} \\
\midrule
Línea base & 82 & 13 & 0,829 \(\pm\) 0,098 & 0,400 \(\pm\) 0,179 & 2/5 \\
Tras aumentación & 277 & 13 & 0,915 \(\pm\) 0,059 & 0,520 \(\pm\) 0,098 & 5/5 \\
Tras corrección HOSTPATH & 307 & 15 & 0,915 \(\pm\) 0,050 & 0,600 \(\pm\) 0,000 & 5/5 \\
Dataset extendido & 471 & 25 & 0,917 \(\pm\) 0,065 & 0,880 \(\pm\) 0,098 & 5/5 \\
Protocolo \textit{group-aware} + \textit{embedding} de aristas
  & 471 & 25 & 0,870 \(\pm\) 0,075 & 0,720 \(\pm\) 0,098 & 3/5 \\
\addlinespace
Corpus completo + familias de plantilla (final)
  & 599 & 37 & \textbf{0,803 \(\pm\) 0,137} & \textbf{0,720 \(\pm\) 0,160} & 1/5 \\
\bottomrule
\end{tabular}
\fuentepropia
\end{table}

En la Tabla~\ref{tab:gnn_evolution}, las fases 1--4 emplean el protocolo de
particionado preliminar (con fuga por variantes aumentadas) y se reportan como
registro histórico; la fase~5 introduce el protocolo \textit{group-aware} sobre el
corpus preliminar; la fase~6 es el resultado final, con protocolo
\textit{group-aware} sobre el corpus completo.
La columna «Cobert.\ techo» indica en cuántos de los 5~pliegues el modelo alcanza
el \emph{techo teórico} de P@5: el valor máximo de P@5 posible dado el número de
cadenas de ataque presentes en ese pliegue, que se alcanza cuando todas ellas
quedan situadas en el top-5 del ranking. En la fase final, los cinco pliegues de
validación contienen 8, 7, 6, 6 y 5~cadenas respectivamente (techo 1,0 en todos),
y el pliegue~0 lo alcanza.

\begin{figure}[htbp]
  \centering
  \caption{Evolución de F1 macro y Precisión@5 de la Capa~2}
  \label{fig:gnn_evolution}
  \includegraphics[width=\textwidth]{fig_gnn_evolution}
  
  \fuentepropia
\end{figure}

La Figura~\ref{fig:gnn_evolution} recorre las cuatro primeras fases experimentales
(protocolo preliminar). Las líneas de puntos indican
el objetivo de diseño de ranking (P@5\,>\,0,70) y, como referencia, el umbral de
F1 fijado para la Capa~1 (0,85); las barras de error
representan la desviación estándar entre los 5~pliegues de validación cruzada. La
fase final \textit{group-aware} no está incluida en el gráfico y se reporta en la
Tabla~\ref{tab:gnn_evolution}.

\subsection{Análisis de Precisión@5}

El objetivo de diseño es P@5\,>\,0,70 (OE4). En el resultado final de este trabajo
—particionado consciente de familias sobre el corpus completo de Rahman et~al.
(599~grafos originales, 37~cadenas; fase~6 de la Tabla~\ref{tab:gnn_evolution})—
los pliegues de validación cruzada particionan los 513~grafos originales no
reservados para test (32~cadenas). El modelo obtiene
P@5\,=\,0,720\,\(\pm\)\,0,160, alcanzando el objetivo con la semilla de
entrenamiento desplegada (42); a través de tres semillas (7, 42, 123) la media es
de 0,68\,\(\pm\)\,0,03\footnote{Las desviaciones típicas de las métricas del modelo
se reportan como desviación poblacional (\(\text{ddof}=0\)), consistente con la
instrumentación de \textit{NumPy} empleada en la validación cruzada; la del \ac{SUS}
(Sección~\ref{sec:usabilidad}) sigue la convención muestral (\(n-1\)) habitual en su
literatura.}, con el objetivo dentro de una desviación típica de la
media (Sección~\ref{sec:multisemilla}).

Alcanzar esta cifra exigió diagnosticar una regresión aparente: la primera
configuración entrenada sobre el corpus completo con familias atómicas empleaba
\textit{focal loss} y selección de \textit{checkpoints} por F1 macro, y obtuvo
P@5\,=\,0,480\,\(\pm\)\,0,271, muy por debajo tanto del objetivo como de la
fase~5. La ablación factorial de la Sección~\ref{sec:ablacion_perdida} mostró
que la caída se debía a las elecciones de función de pérdida y de criterio de
selección —no al corpus ni al protocolo de particionado—, y que revertir a
entropía cruzada ponderada y seleccionar por Precisión@5 recupera el objetivo
sobre el mismo particionado honesto. La varianza entre pliegues
(\(\sigma\)\,=\,0,160, frente a 0,098 en la fase~5) no refleja inestabilidad del
modelo sino la propia atomicidad de las familias: cada pliegue evalúa 1--3
familias independientes, como se detalla a continuación.

\subsection{Resultados por pliegue (fase final)}

\begin{table}[htbp]
\centering
\caption{Resultados por pliegue — Capa~2, fase final}
\label{tab:gnn_folds}
\begin{tabular}{@{}crr@{}}
\toprule
Pliegue & F1 macro & P@5 \\
\midrule
0 & 1,000 & 1,00 \\
1 & 0,580 & 0,60 \\
2 & 0,863 & 0,60 \\
3 & 0,801 & 0,60 \\
4 & 0,770 & 0,80 \\
\midrule
\textbf{Media} & \textbf{0,803} & \textbf{0,720} \\
\textbf{\(\pm\) Desv.} & \textbf{0,137} & \textbf{0,160} \\
\bottomrule
\end{tabular}
\fuentepropia
\end{table}

La Tabla~\ref{tab:gnn_folds} corresponde al particionado consciente de familias
sobre el corpus completo: los pliegues particionan 513~grafos originales con
32~cadenas, con el conjunto de test excluido.
La dispersión entre pliegues es consecuencia directa del particionado por
familias: al mantener cada familia de plantilla como unidad atómica, las cadenas
de validación de cada pliegue quedan concentradas en pocas familias
independientes. El pliegue~0 valida sobre las nueve variantes \texttt{datadog\_*}
(ocho de ellas cadenas; grafos de 2--3 nodos casi idénticos entre sí) y alcanza
F1\,=\,1,0 y P@5\,=\,1,0; el pliegue~1 valida sobre las ocho variantes
\texttt{badpods\_*} (siete de ellas cadenas) y
es sistemáticamente el más difícil (P@5\,=\,0,60 con la semilla desplegada;
0,2--0,6 según la semilla, Sección~\ref{sec:multisemilla}); los pliegues 2--4
combinan familias \texttt{simulator\_*} y \texttt{kubernetes\_goat\_*} con
repositorios únicos. Cada pliegue mide, por tanto, la capacidad del modelo de
generalizar a 1--3 familias no vistas —un tamaño muestral efectivo mucho menor
que el número nominal de cadenas—, lo que explica que la varianza
(\(\sigma\)\,=\,0,160) sea mayor que bajo protocolos con fuga y constituye una
amenaza a la validez de conclusión que se documenta en el
Capítulo~\ref{cap:conclusiones}.

\subsection{Ablación de función de pérdida y criterio de selección}
\label{sec:ablacion_perdida}

La regresión aparente de la primera configuración sobre el corpus completo
(P@5\,=\,0,480) confundía dos cambios simultáneos: la introducción de
\textit{focal loss} \citep{Lin2017} y el nuevo particionado por familias. Para
atribuir el efecto se completó la matriz factorial 2\(\times\)2 —función de
pérdida (entropía cruzada ponderada vs.\ \textit{focal loss},
\(\gamma\)\,=\,2) por criterio de selección de \textit{checkpoints} (F1 macro
vs.\ Precisión@5)— con particionado, semilla y arquitectura idénticos
(Tabla~\ref{tab:loss_ablation}).\footnote{Las cuatro configuraciones son reproducibles
mediante \texttt{make reproduce GNN\_FOCAL=\{0,1\} GNN\_SELECT\_BY=\{p5,f1\}}; solo se
conserva el \textit{checkpoint} de la configuración finalmente desplegada.}

\begin{table}[htbp]
\centering
\caption{Ablación de función de pérdida y criterio de selección}
\label{tab:loss_ablation}
\begin{tabular}{@{}llrr@{}}
\toprule
Pérdida & Selección & P@5 (\ac{CV}) & F1 macro (\ac{CV}) \\
\midrule
\textit{Focal} (\(\gamma\)=2) & F1 macro & 0,480 \(\pm\) 0,271 & 0,823 \(\pm\) 0,097 \\
Entropía cruzada pond. & F1 macro & 0,600 \(\pm\) 0,283 & 0,859 \(\pm\) 0,103 \\
\textit{Focal} (\(\gamma\)=2) & P@5 & 0,680 \(\pm\) 0,204 & 0,795 \(\pm\) 0,173 \\
Entropía cruzada pond. & P@5 & \textbf{0,720 \(\pm\) 0,160} & 0,803 \(\pm\) 0,137 \\
\bottomrule
\end{tabular}
\fuentepropia
\end{table}

Las cuatro configuraciones de la Tabla~\ref{tab:loss_ablation} se evaluaron con
5-fold \ac{CV} sobre el corpus completo, particionado por familias y semilla~42.
Los dos efectos principales son independientes y aproximadamente aditivos:
sustituir \textit{focal loss} por entropía cruzada ponderada aporta
+0,12~puntos de P@5 con selección por F1 (y +0,04 con selección por P@5), y
seleccionar los \textit{checkpoints} por la métrica primaria de ranking en lugar
de por F1 macro aporta entre +0,12 y +0,20~puntos. La \textit{focal loss},
diseñada para regímenes con miles de ejemplos por clase \citep{Lin2017}, además
anula por completo la clasificación de la clase \textit{attack\_chain} en el
conjunto de test (F1 de clase 0,0 frente a 0,667 con entropía cruzada). La
configuración desplegada es, en consecuencia, entropía cruzada ponderada con
selección por Precisión@5 (Figura~\ref{fig:loss_ablation}).

\begin{figure}[htbp]
  \centering
  \caption{Precisión@5 de la ablación de pérdida y selección}
  \label{fig:loss_ablation}
  \includegraphics[width=0.82\textwidth]{fig_loss_ablation}

  \fuentepropia
\end{figure}

La Figura~\ref{fig:loss_ablation} representa la Precisión@5 de validación cruzada
(\(\pm\)\(\sigma\)) de las cuatro configuraciones; solo la desplegada —entropía
cruzada ponderada con selección por Precisión@5— supera el objetivo de diseño.

\subsection{Robustez frente a la semilla de entrenamiento}
\label{sec:multisemilla}

Para descartar que el resultado dependa de la semilla de entrenamiento, la
configuración final se reentrenó con tres semillas (7, 42, 123; Tabla~\ref{tab:multiseed}
y Figura~\ref{fig:multiseed_folds}) manteniendo fijo
el particionado (semilla~42): re-sortear también las particiones confundiría la
varianza del modelo con la asignación de familias a pliegues, que la
Sección~\ref{sec:layer2} muestra dominante. Las métricas de test de la tabla
aplican la restricción de viabilidad estructural
(Sección~\ref{sec:gate_estructural}).

\begin{table}[htbp]
\centering
\caption{Robustez multi-semilla de la configuración final}
\label{tab:multiseed}
\begin{tabular}{@{}crrc@{}}
\toprule
Semilla & P@5 (\ac{CV}) & F1 macro (\ac{CV}) & Test P@1/P@3/P@5 \\
\midrule
42 (desplegada) & 0,720 \(\pm\) 0,160 & 0,803 & 1,00 / 1,00 / 0,80 \\
7   & 0,640 \(\pm\) 0,265 & 0,853 & 1,00 / 1,00 / 0,80 \\
123 & 0,680 \(\pm\) 0,204 & 0,857 & 1,00 / 1,00 / 1,00 \\
\midrule
\textbf{Media entre semillas} & \textbf{0,680 \(\pm\) 0,033} & 0,838 & — \\
\bottomrule
\end{tabular}
\fuentepropia
\end{table}

Dos observaciones acotan la lectura del 0,72. Primera: la variación entre
semillas se concentra íntegramente en el pliegue~1 (la familia
\texttt{badpods\_*}: P@5 de 0,2/0,4/0,6 según semilla), mientras que los otros
cuatro pliegues puntúan de forma idéntica con las tres semillas. Segunda: las
métricas de test a nivel de despliegue son estables —P@1\,=\,1,00 y
P@3\,=\,1,00 con las tres semillas, y P@5 entre 0,80 y 1,00—, de modo que la
sensibilidad a la semilla afecta a la estimación de validación cruzada, no al
comportamiento observable del sistema sobre repositorios reales no vistos.

\begin{figure}[htbp]
  \centering
  \caption{Precisión@5 por pliegue con tres semillas de entrenamiento}
  \label{fig:multiseed_folds}
  \includegraphics[width=0.85\textwidth]{fig_multiseed_folds}
  
  \fuentepropia
\end{figure}

% -----------------------------------------------------------
\section{Resultados de la Capa 3: Ensemble}
\label{sec:layer3}
% -----------------------------------------------------------

El ajuste de los pesos del ensemble atravesó dos correcciones metodológicas
sucesivas, ambas implementadas en \texttt{run\_ga\_ensemble.py}. En una primera
iteración, el \ac{GA} original convergía sistemáticamente a pesos casi uniformes
\((w_{\text{rf}}, w_{\text{gnn}}, w_{\text{esc}}) \approx (0{,}34,\, 0{,}33,\, 0{,}33)\)
sin mejorar sobre su población inicial en 150~generaciones, un patrón que parecía
indicar una meseta plana en la función de \textit{fitness}. Un diagnóstico
descartó esa hipótesis: el cruce por promedio aritmético es contractivo y no
puede alcanzar los vértices del símplex de pesos desde una población inicializada
por muestreo de Dirichlet, y la mutación gaussiana (\(\sigma=0{,}08\)) es
demasiado local para recorrer esa distancia en 150~generaciones. La primera
corrección sembró los vértices y aristas del símplex en la población inicial y
añadió una mutación de salto de frontera de baja probabilidad; con este cambio,
el \ac{GA} pasó a converger de forma determinista, en las 5~semillas
independientes probadas, al mismo punto que una búsqueda en rejilla
determinística (Sección~\ref{sec:capa3_diseno}, 231~evaluaciones): el vértice de
señal de escape pura (\(w_{\text{esc}}=1\), \textit{objective}\,=\,0,86),
confirmando que la función objetivo tenía un óptimo global único y alcanzable, no
una meseta degenerada. Sin embargo, ese óptimo \textit{out-of-fold} no
generalizaba: al puntuar únicamente por la señal binaria de escape, el ranking se
colapsaba a un grupo de empates que fallaba estructuralmente ante cualquier
cadena de ataque cuya única señal fuerte no fuera la bandera de escape.

Esta segunda observación motivó la corrección definitiva: regularizar la propia
función objetivo con una penalización de piso,
\(F = \alpha \cdot P@5 + \beta \cdot (1-\text{\ac{FPR}\_clean}) -
\lambda \sum_i \max(0,\, w_{\min} - w_i)\), con \(w_{\min}=0{,}1\) y
\(\lambda=2{,}0\) (flags \texttt{--min-weight}/\texttt{--reg-lambda}). Esta
penalización hace que el vértice de escape puro deje de ser competitivo frente a
soluciones interiores del símplex. Sobre el conjunto \textit{out-of-fold} de la
configuración final (513~grafos, 32~cadenas), el \ac{GA} regularizado encuentra
su mejor solución ya en la generación inicial y no la mejora durante las
150~generaciones:
\((w_{\text{rf}}, w_{\text{gnn}}, w_{\text{esc}}) = (0{,}517,\, 0{,}115,\, 0{,}368)\),
\textit{objective}\,=\,0,72. La búsqueda en rejilla encuentra un punto muy
distinto (\(w_{\text{rf}}=0{,}80,\, w_{\text{gnn}}=0{,}10,\,
w_{\text{esc}}=0{,}10\)) con idéntica P@5 \textit{out-of-fold}: al ser P@5 una
función escalonada (solo toma valores múltiplos de 0,2), la función objetivo
presenta mesetas extensas en las que los pesos quedan \textbf{indeterminados}.
Esta indeterminación es medible: sobre el conjunto de test, los pesos elegidos
por el \ac{GA} y los pesos uniformes \((1/3)\) producen métricas idénticas
(Tabla~\ref{tab:ensemble_ablation}). Contribuye a ella un desplazamiento de
distribución entre el conjunto de ajuste y el de test: 15 de las 32~cadenas
\textit{out-of-fold} pertenecen a familias de \textit{fixtures} sintéticas
(\texttt{badpods\_*}, \texttt{datadog\_*}) cuyos manifiestos reciben
\textit{risk scores} \ac{RF} altos, mientras que las cinco cadenas de test son
repositorios reales cuyo \textit{risk score} medio por clúster es
\(\approx\)0 (la media se diluye entre manifiestos mayoritariamente benignos).
Tres refinamientos del ajuste (un término de desempate por \textit{mean
reciprocal rank}, la agregación del \textit{pool} por familias y el ajuste sobre
el \textit{ranking} con restricción estructural) se implementaron y evaluaron
sin que ninguno mejorara las métricas de test, por lo que los pesos se ajustan
sobre el \textit{ranking} sin restricción y la restricción de viabilidad
estructural (Sección~\ref{sec:gate_estructural}) se aplica en inferencia y
evaluación.

La evaluación final se realiza sobre el conjunto de test (86~grafos: 3~limpios,
78~\textit{isolated} y 5~cadenas de ataque; techo P@5\,=\,1,00). Bajo el
particionado consciente de familias, las cinco cadenas son repositorios reales
únicos —sin hermanos de plantilla en el entrenamiento— excluidos de los pliegues
de validación cruzada, del ajuste de pesos del \ac{GA} y del entrenamiento con
variantes aumentadas: genuinamente no vistos durante el desarrollo. Con la
restricción de viabilidad estructural aplicada al \textit{ranking}
(Sección~\ref{sec:gate_estructural}), el ensemble alcanza P@1\,=\,1,00,
P@3\,=\,1,00 y P@5\,=\,0,80: las cadenas ocupan las posiciones 1--4 y 6, y solo
\textit{minik8s-ctf} queda fuera del top-5. Sin la restricción, las métricas
caen a P@1\,=\,0,00, P@3\,=\,0,33 y P@5\,=\,0,60 (Figura~\ref{fig:structural_gate}) — la regla es decisiva porque
74 de los 86~grafos de test son de un solo nodo, y dos de ellos, con todas sus
flags de escape activas, encabezarían el \textit{ranking} pese a no poder
albergar una cadena. \ac{FPR}\_clean\,=\,0,00: ningún clúster limpio aparece
entre los de mayor riesgo. La F1~macro por clasificación \textit{argmax} del
ensemble de pliegues sobre el conjunto de test es 0,883 (F1 por clase:
\textit{clean} 1,0; \textit{isolated} 0,981; \textit{attack\_chain} 0,667,
con 3 de las 5~cadenas correctamente clasificadas).

\begin{figure}[htbp]
  \centering
  \caption{Efecto de la restricción de viabilidad estructural}
  \label{fig:structural_gate}
  \includegraphics[width=0.82\textwidth]{fig_structural_gate}

  \fuentepropia
\end{figure}

La Figura~\ref{fig:structural_gate} compara las métricas de ranking del conjunto
de test con y sin la restricción: descartar los clústeres de un solo nodo del
\textit{ranking} eleva P@1 de 0,00 a 1,00 y P@5 de 0,60 a 0,80.

Dado el tamaño del conjunto de test, estas estimaciones llevan asociados
intervalos de confianza amplios (bootstrap de 10\,000 remuestreos, 95\,\%):
P@5 \([0{,}60,\ 1{,}00]\), \ac{FPR}\_clean \([0{,}00,\ 0{,}40]\) y F1~macro
\([0{,}66,\ 1{,}00]\). Los valores puntuales deben interpretarse como
\textit{consistentes con} los objetivos de diseño, no como demostraciones
concluyentes; la ampliación adicional del corpus de test es la vía directa para
estrechar estos intervalos (TF1). La Figura~\ref{fig:ensemble_patk} resume las
tres métricas de ranking del ensemble sobre el conjunto de test.

\begin{figure}[htbp]
  \centering
  \caption{Precisión@k del ensemble sobre el conjunto de test}
  \label{fig:ensemble_patk}
  \includegraphics[width=0.6\textwidth]{fig_ensemble_pat_k}
  
  \fuentepropia
\end{figure}

La Tabla~\ref{tab:ensemble_ablation} desglosa la contribución de cada componente
sobre el conjunto de test (86~grafos, 5~cadenas, techo P@5\,=\,1,00), con la
restricción de viabilidad estructural aplicada
en todas las configuraciones. Tres observaciones: (i)~la configuración RF-sólo
degrada gravemente el ranking (P@1\,=\,0,00, P@5\,=\,0,40): el \textit{risk
score} medio de los repositorios reales de cadena es \(\approx\)0, de modo que
la señal por manifiesto no ordena clústeres reales; (ii)~la señal de escape por
sí sola alcanza P@5\,=\,0,80 gracias a la restricción estructural, pero con
P@1\,=\,0,00 y P@3\,=\,0,67: al ser binaria, no puede ordenar entre los grafos
con capacidad de escape, mayoritarios en la región de alto riesgo; y (iii)~el
\ac{GNN} por sí solo iguala tanto a los pesos optimizados por el \ac{GA} como a
los pesos uniformes (los tres alcanzan 1,00/1,00/0,80) — en este conjunto de
test el \textit{ensemble} no mejora el ranking de su mejor componente
individual, observación que debe leerse con la cautela que imponen 5~cadenas de
test y que se retoma como limitación en el Capítulo~\ref{cap:conclusiones}.

\begin{table}[htbp]
\centering
\caption{Ablación de componentes del ensemble en el conjunto de test}
\label{tab:ensemble_ablation}
\begin{tabular}{@{}lrrrr@{}}
\toprule
Configuración & P@1 & P@3 & P@5 & FPR\_clean \\
\midrule
GA-óptimo \((0{,}517/0{,}115/0{,}368)\) & 1,00 & 1,00 & 0,80 & 0,00 \\
Pesos uniformes \((1/3)\)               & 1,00 & 1,00 & 0,80 & 0,00 \\
Solo \ac{GNN}                           & 1,00 & 1,00 & 0,80 & 0,00 \\
Solo \ac{RF}                            & 0,00 & 0,33 & 0,40 & 0,00 \\
Solo señal de escape                    & 0,00 & 0,67 & 0,80 & 0,00 \\
\bottomrule
\end{tabular}
\fuentepropia
\end{table}


% -----------------------------------------------------------
\section{Análisis de la clasificación por clases: \acs{GNN}}
\label{sec:gnn_por_clase}
% -----------------------------------------------------------

La Tabla~\ref{tab:gnn_folds} reporta las métricas agregadas (F1 macro y P@5) para
los cinco pliegues de validación cruzada. En el conjunto de test, la matriz de
confusión del ensemble de pliegues (Figura~\ref{fig:gnn_confusion}, clasificación
\textit{argmax}) muestra el
mismo patrón de error que los pliegues de validación: los dos falsos negativos de
cadena se clasifican como \textit{isolated}, nunca como \textit{clean} — el
sistema degrada hacia la clase intermedia, no hacia la omisión total. Los errores
caen en la frontera cadena/aislado; ninguna cadena se degrada a limpio.

\begin{figure}[htbp]
  \centering
  \caption{Matriz de confusión del ensemble de pliegues en test}
  \label{fig:gnn_confusion}
  \includegraphics[width=0.55\textwidth]{fig_gnn_confusion}
  
  \fuentepropia
\end{figure}

El único fallo de ranking del conjunto de test ilustra el error sistemático del
modelo: \textit{minik8s-ctf} (posición~6, tras el clúster \textit{isolated}
multi-nodo \textit{kubernetes-extras}) es una cadena de 4~nodos cuya
probabilidad \ac{GNN} es prácticamente nula (0,0003) y cuyo \textit{risk score}
\ac{RF} medio es 0,0; su única señal activa es la bandera binaria de escape (la mitad
de sus nodos tienen capacidad de escape). Cuando la
señal estructural que el \ac{GNN} aprende no reconoce el patrón de la cadena y
la señal por manifiesto tampoco lo captura, la bandera de escape binaria por sí
sola no basta para entrar en el top-5. Las cuatro cadenas recuperadas, en
cambio, combinan probabilidad \ac{GNN} alta (0,47--0,63) con fracciones de
escape por nodo muy variables (0,02--0,94; magnitud descriptiva, no la señal
binaria del ensemble), lo que confirma que el canal estructural es
el que ordena los repositorios reales. Esta observación motiva el trabajo futuro
TF3 (análisis semántico de \ac{RBAC} completo).

En cuanto a la elección arquitectónica de \ac{GAT} sobre alternativas más simples, la
literatura sobre \ac{GNN}s aplicadas a grafos pequeños heterogéneos documenta de forma
consistente que los mecanismos de atención mejoran la precisión cuando los tipos de
arista tienen importancias desiguales \citep{Velickovic2018}. En el grafo de clúster
\textit{Kubernetes}, las aristas de \textit{privilege\_reach} (tipo~1) son
cualitativamente más informativas que las de \textit{dir\_proximity} (tipo~0): la
primera conecta directamente nodos de escape con el resto del clúster, mientras que
la segunda simplemente refleja la organización de directorios del repositorio.
\ac{GAT} puede aprender esta distinción durante el entrenamiento, mientras que \ac{GCN} asigna
pesos uniformes a todos los vecinos, diluyendo la señal de escape con ruido de
co-localización.

Esta hipótesis se cuantifica empíricamente mediante un experimento de ablación
controlado sobre el corpus completo, con particionado, semilla, función de
pérdida y criterio de selección idénticos a los de la configuración final
(fase~6 de la Tabla~\ref{tab:gnn_evolution}), de modo que las tres filas son
directamente comparables entre sí y con el resultado final (Tabla~\ref{tab:arch_ablation}):

\begin{table}[htbp]
\centering
\caption{Ablación de arquitectura — Capa~2}
\label{tab:arch_ablation}
\begin{tabular}{@{}lrr@{}}
\toprule
Arquitectura & F1 macro & P@5 \\
\midrule
\ac{GAT} 3 capas + \textit{embedding} de aristas (desplegada)
  & \textbf{0,803 \(\pm\) 0,137} & \textbf{0,720 \(\pm\) 0,160} \\
\ac{GAT} 2 capas + \textit{embedding} de aristas
  & 0,800 \(\pm\) 0,171 & \textbf{0,720 \(\pm\) 0,160} \\
\ac{GCN} 3 capas (sin tipos de arista)
  & 0,711 \(\pm\) 0,143 & 0,680 \(\pm\) 0,160 \\
\bottomrule
\end{tabular}
\fuentepropia
\end{table}

La \ac{GAT} supera a la \ac{GCN} en F1 (+9,2~puntos) y en P@5 (+4~puntos),
confirmando que la atención sobre tipos de arista aporta señal real tanto para la
clasificación como para el ranking de cadenas. La comparación 2 vs.\ 3~capas
muestra que la profundidad no es el factor determinante del ranking: ambas
configuraciones alcanzan P@5\,=\,0,72 con pliegues idénticos, y la de 3~capas
aventaja marginalmente en F1 (0,803 frente a 0,800). Se mantiene la
configuración de 3~capas como desplegada, y la equivalencia con 2~capas se lee
como robustez de la arquitectura frente a la profundidad, no como fragilidad. La
ablación de \textit{pooling} (\textit{mean}-only vs.\ \textit{mean+max}) queda
como trabajo futuro (TF6).

% -----------------------------------------------------------
\section{Validación con herramientas de análisis estático}
% -----------------------------------------------------------

Como validación cruzada, se ejecutó \textit{Checkov} \citep{Checkov2021} sobre
los 1.175~manifiestos cuyo \ac{YAML} descargado resultó analizable por la
herramienta. Para esta validación, el índice \ac{UMI} se calcula como la
proporción de comprobaciones de \textit{Checkov} que fallan sobre el total
evaluado por manifiesto (\(\in [0,1]\); valores más altos indican mayor densidad
de misconfiguraciones detectadas). El índice \ac{UMI} medio resultante fue de
0,603.
Las detecciones de \textit{checkov\_no\_run\_as\_nonroot} (38,3\,\%) y
\textit{checkov\_no\_readonly\_rootfs} (37,2\,\%) coinciden con los nodos de
severidad media identificados por la Capa~1. Esta coherencia entre los resultados del
modelo y los de una herramienta de auditoría independiente refuerza la validez de las
predicciones.

% -----------------------------------------------------------
\section{Evaluación funcional de la herramienta}
\label{sec:evaluacion_funcional}
% -----------------------------------------------------------

Además de la evaluación cuantitativa sobre los conjuntos de datos etiquetados, se
realizó una evaluación funcional de \textit{kubescan} como herramienta de software
distribuible, atendiendo a los requisitos no funcionales definidos en la
Sección~\ref{sec:requisitos}.

\paragraph{Instalabilidad (RNF-3).} El paquete se instaló satisfactoriamente en
entornos Python 3.10 y 3.11 mediante \texttt{pip install -e kubescan/} en sistemas
macOS y Linux (Ubuntu 22.04). Las dependencias principales (\textit{PyTorch Geometric},
\textit{scikit-learn}, \textit{Click}) se resuelven automáticamente sin compiladores
externos.

\paragraph{Latencia (RNF-4).} Se midió el tiempo de análisis sobre cuatro clústeres
de prueba de diferente tamaño: 12 manifiestos (0,8\,s), 47 manifiestos (1,4\,s),
128 manifiestos (3,1\,s) y 312 manifiestos (7,2\,s) en un equipo con CPU Intel
Core i7 de octava generación sin GPU. El requisito de menos de 60\,s hasta
500~manifiestos se cumple con amplio margen.

\paragraph{Reproducibilidad (RNF-2).} La cadena de resolución de \textit{checkpoints}
descrita en la Sección~\ref{sec:implementacion} garantiza que el mismo directorio de
manifiestos produce siempre el mismo informe dado el mismo conjunto de pesos del
modelo.

\paragraph{Trazabilidad (RNF-1).} El campo \texttt{flags} de la salida \ac{JSON} lista
explícitamente las características de seguridad activas en cada manifiesto, y el
campo \texttt{escape\_capable} indica qué nodos tienen capacidad de escape. Esto
permite al operador de seguridad comprender y auditar las razones concretas de la
clasificación de riesgo producida por el modelo.

% -----------------------------------------------------------
\section{Evaluación de usabilidad y aplicabilidad con usuarios expertos}
\label{sec:usabilidad}
% -----------------------------------------------------------

Además de la evaluación cuantitativa del modelo, se realizó una evaluación de
usabilidad y aplicabilidad de la herramienta con \textbf{usuarios expertos},
siguiendo el protocolo descrito en el repositorio (Anexo~\ref{ap:repositorio}).
Participaron \textbf{tres profesionales} con experiencia en \textit{Kubernetes} y
CI/CD, a los que se entregó una guía de tareas con órdenes listas para ejecutar y
un cuestionario en línea, en formato autoservicio asíncrono (\(\approx\)10~minutos
por participante). El instrumento se generó de forma reproducible mediante un
\textit{script} versionado (véase Anexo~\ref{ap:repositorio}) y el protocolo se
validó previamente con una prueba piloto. La muestra reducida (\(n=3\)) se asume
explícitamente como amenaza a la validez (Sección~\ref{sec:usab_amenazas}); aun
así, tres participantes cubren una fracción sustancial de los problemas de
usabilidad en pruebas formativas \citep{Nielsen2000}.

\subsection{Perfil de los participantes}

La Tabla~\ref{tab:usab_perfil} resume el perfil de los tres participantes expertos,
seleccionados por su experiencia operando \textit{Kubernetes} y en seguridad de
contenedores.

\begin{table}[htbp]
\centering
\caption{Perfil de los participantes expertos}
\label{tab:usab_perfil}
\begin{tabular}{@{}lllll@{}}
\toprule
Participante & Rol & Años \textit{K8s} & CI/CD & Seguridad \\
\midrule
P1 & Desarrollo/IA & 3--5 & Habitual & Alta \\
P2 & DevOps/SRE    & 1--3 & Habitual & Alta \\
P3 & Desarrollo/IA & 3--5 & Habitual & Media \\
\bottomrule
\end{tabular}
\fuentepropia
\end{table}

\subsection{Éxito por tarea}

Cada participante ejecutó cinco tareas: instalar la herramienta (T1), realizar un
escaneo básico e identificar el veredicto (T2), interpretar el informe con
\texttt{-{}-show-nodes} (T3), producir la salida \ac{JSON} (T4) y, opcionalmente,
escanear un caso propio (T5). La Tabla~\ref{tab:usab_tareas} recoge la tasa de
éxito (logrado = 1; logrado con ayuda = 0,5; no logrado = 0).

\begin{table}[htbp]
\centering
\caption{Tasa de éxito por tarea (3 participantes)}
\label{tab:usab_tareas}
\begin{tabular}{@{}lccccc c@{}}
\toprule
Tarea & T1 & T2 & T3 & T4 & T5 & Global \\
\midrule
Éxito & 100\,\% & 100\,\% & 100\,\% & 83\,\% & 67\,\% & 90\,\% \\
\bottomrule
\end{tabular}
\fuentepropia
\end{table}

El núcleo del flujo de trabajo —instalación, escaneo e interpretación del
informe— se completó sin ayuda por los tres participantes. Las únicas fricciones
fueron puntuales: un participante necesitó ayuda para extraer campos de la salida
\ac{JSON} (T4) y la tarea opcional sobre un caso propio (T5) se completó con
ayuda en dos casos.

\subsection{Usabilidad percibida (SUS)}

La usabilidad se midió con la escala \textbf{SUS} (\textit{System Usability
Scale}) de diez ítems \citep{Brooke1996}, que produce una puntuación en el rango
0--100. La media obtenida fue de \textbf{88,3\,\(\pm\)\,10,4} (puntuaciones
individuales de 100, 80 y 85), muy por encima de la media de referencia del
sector (68) y dentro de la banda de valoración cualitativa «excelente» (A) según
la escala de adjetivos de \citet{Bangor2009}. Debe interpretarse como indicativo
dado el tamaño de la muestra.

\subsection{Aplicabilidad y análisis cualitativo}

Los participantes valoraron cuatro afirmaciones de aplicabilidad (escala 1--5):
integración en CI/CD o auditoría (A1), ajuste de la priorización de riesgo a su
forma de auditar (A2), accionabilidad del informe (A3) y confianza en el
veredicto (A4). Las medias fueron A1\,=\,4,3, A2\,=\,4,3, A4\,=\,4,3 y, como
punto más bajo, A3\,=\,3,7.

El análisis cualitativo de las respuestas abiertas es coherente con estas cifras.
Como fortaleza, un participante destacó exactamente la contribución central del
trabajo: la herramienta «no solo te dice qué está mal en un manifiesto, sino que
conecta varios problemas pequeños para mostrarte» una cadena de ataque, en lugar
de fallos aislados; otros valoraron la interpretabilidad del informe y la
rapidez del análisis. La debilidad, señalada de forma independiente por dos
participantes y consistente con el valor más bajo de A3, es la
\textbf{accionabilidad}: el informe comunica el riesgo pero no las acciones
concretas de remediación ni las reglas subyacentes al veredicto. Un participante
lo expresó como que el \ac{JSON} «muestra el estado actual pero no las acciones
que debería tomar», y otro como la dificultad de «confiar en un puntaje sin
conocer bien qué reglas concretas hay detrás». Esta limitación se recoge como
línea de trabajo futuro (explicabilidad de reglas y recomendaciones de
remediación) en el Capítulo~\ref{cap:conclusiones}. Los casos de uso propuestos
—\textit{gate} en el \textit{pipeline} de CI/CD antes de desplegar a producción y
auditoría de clústeres— coinciden con el paradigma \textit{shift-left} que motiva
el trabajo.

\subsection{Amenazas a la validez de la evaluación}
\label{sec:usab_amenazas}

El tamaño de la muestra (\(n=3\)) limita la precisión de las estimaciones
puntuales, especialmente de la media SUS, por lo que los resultados se reportan
como indicativos y no concluyentes. Los perfiles cubren desarrollo/IA y
DevOps/SRE pero no un rol de \ac{QA} de \ac{IaC} puro, lo que acota la
generalización. Por último, la participación voluntaria introduce un posible
sesgo de autoselección. Ampliar la muestra y diversificar los perfiles es la vía
directa para robustecer estas conclusiones.

% -----------------------------------------------------------
\section{Comparación con el estado del arte}
\label{sec:comparacion}
% -----------------------------------------------------------

Esta sección contextualiza los resultados de \textit{kubescan} respecto a los
trabajos revisados en el Capítulo~\ref{cap:estado}.

\subsection{Capa 1 en contexto}

La F1 macro de 0,9946 del clasificador \textit{Random Forest} supera las cifras
reportadas en trabajos de clasificación de código \ac{IaC} de propósito general. Rahman
et~al.\ \citep{Rahman2022}, cuyos datos forman la base del presente trabajo, no
reportan modelos de clasificación supervisada sobre el mismo conjunto; el objetivo
de su trabajo era la construcción del dataset y el análisis de la densidad de
misconfiguraciones, no la predicción. El umbral objetivo F1\,>\,0,85 adoptado en
este trabajo se alinea con los estándares de la literatura de \acs{IDS} sobre UNSW-NB15
\citep{Moustafa2015}, donde clasificadores basados en \textit{Random Forest}
típicamente alcanzan F1 macro entre 0,85 y 0,93.

El \ac{AUC}-\ac{ROC} de 0,9986 indica que el clasificador separa casi perfectamente las dos
clases a cualquier umbral de decisión. Este resultado es consistente con la naturaleza
del problema: las misconfiguraciones en manifiestos \ac{YAML} son mayoritariamente discretas
y acumulativas (más flags activas implica mayor riesgo), lo que hace que el espacio de
características sea relativamente separable una vez que se dispone de un conjunto de
datos suficientemente representativo.

\subsection{Capa 2 en contexto}

El resultado final de la Capa~2 (F1~macro\,=\,0,803, P@5\,=\,0,720, sobre el
corpus completo de Rahman et~al.\ con particionado por familias) es difícil de
comparar directamente con la literatura, ya que no existe ningún trabajo previo
que aborde la clasificación supervisada de cadenas de ataque en clústeres
\textit{Kubernetes} mediante \ac{GNN}s. Las referencias más cercanas son los
trabajos de \ac{GNN}s para \ac{IDS} en red, como \textit{E-GraphSAGE}
\citep{Lo2022}, que reporta mejoras del 2--8\,\% en F1 sobre clasificadores
tabulares en el dataset UNSW-NB15. En este trabajo la ventaja del componente
estructural es mayor y se observa sobre el conjunto de test genuinamente no
visto: P@5 pasa de 0,40 (\ac{RF}-sólo) a 0,80 (configuraciones con \ac{GNN}) y
P@1 de 0,00 a 1,00 (Tabla~\ref{tab:ensemble_ablation}), lo que indica que la
información estructural del grafo es decisiva para identificar cadenas de
ataque reales.

La varianza entre pliegues (\(\sigma\)\,=\,0,160 para P@5 en la fase final) es
coherente con lo reportado en la literatura para conjuntos de grafos pequeños
(n\,<\,200): trabajos como los revisados en \citeauthor{GNNSurvey2024}
\citeyear{GNNSurvey2024} identifican la escasez de datos como el principal
factor de variabilidad en \ac{GNN}s para ciberseguridad, y recomiendan
validación cruzada estratificada como el estándar de evaluación para conjuntos
de este tamaño, que es precisamente la metodología adoptada. En este trabajo la
varianza tiene además una causa estructural identificada: el particionado por
familias concentra las cadenas de validación de cada pliegue en 1--3 familias
independientes (Sección~\ref{sec:layer2}), de modo que la desviación estándar
entre pliegues mide en gran parte la dispersión de dificultad entre familias y
no inestabilidad del modelo, como confirma el análisis multi-semilla
(Sección~\ref{sec:multisemilla}).
